\documentclass[10pt,a4paper]{article}
\usepackage[T1]{fontenc}
\usepackage[french]{babel}
\usepackage{lmodern,amsmath,amssymb,booktabs,graphicx,microtype}
\usepackage[margin=15mm]{geometry}
\pagestyle{empty}
\setlength{\parindent}{0pt}
\setlength{\parskip}{1pt}
\setlength{\topsep}{4pt}
\setlength{\abovedisplayskip}{4pt}
\setlength{\belowdisplayskip}{4pt}
\newcommand{\sectionline}[1]{\vspace{4pt}\textbf{#1}\par}
\begin{document}
{\Large\bfseries Audit GARCH}\hfill {\small\bfseries Yosri BEN HALIMA}\par
{\small Asymétrie de la volatilité du S\&P 500, rendements hebdomadaires}\par
\vspace{3pt}\hrule height 0.35pt\vspace{5pt}

J'audite notre GARCH(1,1) à moyenne constante sur les rendements logarithmiques
hebdomadaires du S\&P 500 en pourcentage et le compare au module Python \texttt{arch}.
J'examine ensuite l'asymétrie de la volatilité avec le GJR-GARCH et les courbes d'impact des chocs.

\sectionline{1. Audit}
\begin{center}
\small
\renewcommand{\arraystretch}{1.08}
\textbf{Tableau 1 : audit des six points de contrôle}\par
\begin{tabular}{@{}p{0.18\linewidth}p{0.09\linewidth}p{0.67\linewidth}@{}}
\toprule
Point & Verdict & Vérification et correction \\
\midrule
Résidus & Correct & La récurrence utilise $e_t=r_t-\mu$. \\
Variance initiale & Correct & La variance empirique initialise la récurrence, avec le diviseur $T$. Ici, $h_1=5{,}2774$. \\
Vraisemblance & Correct & Minimisation de l'opposé de la somme des log-vraisemblances gaussiennes. \\
Contraintes & Correct & Transformations exponentielle/softmax : $\omega>0$, $\alpha,\beta>0$, $\alpha+\beta<1$. Les zéros exacts sont exclus. \\
Erreurs-types & Erroné & La Hessienne numérique a été remplacée par BHHH : inverse de la somme des produits extérieurs des scores, moyenne incluse, puis méthode delta. \\
Échelle & Correct & Facteur 100 conservé après le passage aux rendements logarithmiques. \\
\bottomrule
\end{tabular}
\par\vspace{2pt}{\footnotesize Les verdicts portent sur le code initial. La comparaison utilise la version corrigée.}
\end{center}

\vspace{3pt}
{\centering\small
\textbf{Tableau 2 : estimations GARCH(1,1), notre implémentation et le module \texttt{arch}}\par
\begin{tabular}{lrrrrr}
\toprule
 & Notre implémentation & É.-t. BHHH & GARCH \texttt{arch} & Écart absolu & Écart relatif (\%) \\
\midrule
$\mu$    & $0{,}260789$ & $0{,}035543$ & $0{,}260804$ & $0{,}000015$ & $0{,}01$ \\
$\omega$ & $0{,}320175$ & $0{,}048116$ & $0{,}320190$ & $0{,}000015$ & $0{,}00$ \\
$\alpha$ & $0{,}211234$ & $0{,}015736$ & $0{,}211268$ & $0{,}000034$ & $0{,}02$ \\
$\beta$  & $0{,}739121$ & $0{,}018465$ & $0{,}739098$ & $0{,}000022$ & $0{,}00$ \\
\bottomrule
\end{tabular}
\par}\vspace{3pt}

Les coefficients concordent à deux décimales. Les log-vraisemblances valent
$-5084{,}787$ et $-5084{,}785$. Pour $\omega$, l'écart entre erreurs-types BHHH
($0{,}0481$) et robuste ($0{,}0845$) invite à la prudence sur l'inférence.

ARCH-LM rejette la variance constante ($12$ retards, $LM=374{,}937$,
$p=7{,}60\times10^{-73}$), puis ne détecte plus d'effet ARCH après filtrage ($p=0{,}712$).
Ces contrôles confortent le code corrigé. La loi gaussienne et les prévisions restent à tester.

\sectionline{2. Asymétrie}
Pour le GJR, $\widehat\gamma=0{,}272354$, avec une erreur-type robuste de $0{,}056559$,
$t=4{,}815$ et $p=1{,}47\times10^{-6}$ (asymptotique). La corrélation entre le rendement
de la semaine et la valeur absolue du rendement suivant vaut $-0{,}2243$ (2\,386 paires).
Cette corrélation négative suggère que les baisses tendent à précéder des mouvements
plus amples la semaine suivante, sans établir de causalité.

Pour isoler l'effet du choc, je fixe la variance à son niveau moyen de long terme
(variance inconditionnelle) : $\bar h_G=6{,}449$ et $\bar h_J=5{,}360$ (pourcentage au carré).
\begin{center}
\includegraphics[width=\linewidth]{news_impact.pdf}
\end{center}
À choc de même ampleur, le GJR réagit davantage du côté négatif. Le GARCH
reste symétrique. Les courbes diffèrent aussi du côté positif puisque chaque
estimation a ses coefficients et sa variance de long terme. Celle du GJR suppose des innovations symétriques.

\sectionline{3. Erreurs de l'IA et corrections}
L'IA avait commis deux erreurs, corrigées ensuite :
\begin{itemize}
\setlength{\itemsep}{0pt}
\setlength{\parsep}{0pt}
\item des rendements simples au lieu de rendements logarithmiques.
\item une Hessienne numérique au lieu de la méthode BHHH pour les erreurs-types.
\end{itemize}
\end{document}
